%% ── 1-Р БҮЛЭГ: УДИРТГАЛ ─────────────────────────────────────
\chapter{Удиртгал}
\label{ch:intro}

\section{Үндэслэл, ач холбогдол}

Өдөр тутмын мэдээлэл, хийх ажлын жагсаалт, төлөвлөгөөгөө хурдан тэмдэглэж хаанаас ч
хандах хэрэгцээ нэмэгдэж байгаа ч одоогийн дижитал тэмдэглэлийн аппликейшнүүд
хэрэглэгчийн бодит шаардлагад нийцсэн ухаалаг функц дутмаг байна. Тиймээс хиймэл оюун
ашиглан тэмдэглэлийн утга агуулгад тохирсон стикер автоматаар үүсгэж, тэдгээр стикерийг
календарь дээр тухайн өдөрт нь байрлуулан, автоматаар мэдэгдэл илгээдэг шинэ төрлийн
ухаалаг тэмдэглэлийн аппликейшн хөгжүүлэх шаардлага бий болж байна.

Жишээлбэл, хэрэглэгч ``3 сарын 8-нд Sunday уулзалт байна'' гэж тэмдэглэл оруулахад систем
тухайн агуулгад тохирсон стикерийг хиймэл оюун ашиглан автоматаар үүсгэж, тэр стикерийг
календарийн 3 сарын 8-ны нүдэнд байрлуулна. Тохируулсан цагт нь push notification
илгээх эсэхийг систем өөрөө шийдэж, хэрэглэгчид сануулга хүргэнэ. Энэ нь хэрэглэгчийн
мэдээллийг илүү ойлгомжтой, хурдан сэргээн санах боломжтой болгоно.

\section{Судалгааны асуудал}

Одоогийн дижитал тэмдэглэлийн системүүд нь тэмдэглэл хадгалах, засах, хайх үндсэн
функцтэй ч дараах асуудлуудтай байна:
\begin{itemize}
  \item Хэрэглэгчийн сэтгэл хөдлөл, mood-ийг харгалзан үздэггүй
  \item Тэмдэглэлийн агуулгад тохирсон визуал стикер автоматаар үүсдэггүй
  \item Үүсгэсэн стикерийг календарийн тодорхой өдөрт холбох боломжгүй
  \item Тэмдэглэлийн агуулгад суурилсан ухаалаг сануулгын систем байдаггүй
  \item Монгол хэрэглэгчдэд тохирсон орон нутгийн шийдэл дутмаг
\end{itemize}

\section{Зорилго, зорилт}

Энэхүү судалгааны ажлын үндсэн зорилго нь хиймэл оюун ухааны технологийг ашиглан
хэрэглэгчийн тэмдэглэлийн агуулга болон mood-д нийцсэн стикер дизайн автоматаар
үүсгэж, тэдгээр стикерийг \textbf{календарийн тодорхой өдөрт} байрлуулан, тухайн
өдрийн сануулга, мэдэгдлийн системийг систем өөрөө удирдаж ажиллуулдаг дижитал
ухаалаг тэмдэглэлийн дэвтрийн аппликейшныг боловсруулахад оршино.

Судалгааны зорилтуудыг дараах байдлаар тодорхойлов:

\begin{enumerate}
  \item Дижитал тэмдэглэлийн аппликейшний хэрэглэгчийн хэрэгцээ, шаардлагыг судалж тодорхойлох.
  \item Flutter (Dart) ашиглан хэрэглэгчдэд ойлгомжтой, хялбар UI/UX загвар боловсруулах.
  \item Тэмдэглэл үүсгэх, засварлах, устгах, харах үндсэн CRUD үйлдлүүдийг хэрэгжүүлэх.
  \item Хэрэглэгчийн сэтгэлийн байдал (mood) дээр үндэслэн тэмдэглэлийг ангилах.
  \item \textbf{Хиймэл оюун ашиглан тэмдэглэлийн агуулгад тохирсон стикер дизайн автомат үүсгэх.}
  \item \textbf{Стикерийг календарийн тодорхой өдөрт холбож харуулах calendar view хэрэгжүүлэх.}
  \item \textbf{Тухайн өдрийн тэмдэглэлд суурилсан push notification автоматаар илгээх систем нэвтрүүлэх.}
  \item Node.js болон MongoDB ашиглан өгөгдлийг хадгалах болон дамжуулах серверийн шийдэл нэвтрүүлэх.
  \item JWT болон bcrypt ашиглан хэрэглэгчийн өгөгдлийн аюулгүй байдал, нууцлалыг хангах.
  \item Системийн гүйцэтгэлийг тестлэж, алдааг засварлан оновчтой болгох.
\end{enumerate}

\section{Судалгааны объект, хамрах хүрээ}

Судалгааны объект нь Шинэ Монгол Технологийн Коллежийн оюутнуудын өдөр тутмын
тэмдэглэл хөтлөлтийн хэрэгцээ юм. Системийн туршилтыг коллежийн оюутнуудад
гаргаж, хэрэглэгчийн санал хүсэлтийг цуглуулан үнэлгээ хийв.

\textbf{Хамрах хүрээ:}
\begin{itemize}
  \item Мобайл аппликейшн (Android болон iOS)
  \item REST API backend систем
  \item MongoDB үүлэн өгөгдлийн сан
  \item Хэрэглэгчийн баталгаажуулалт болон аюулгүй байдал
  \item Тэмдэглэлийн удирдлага (CRUD) болон mood tracking
  \item \textbf{AI sticker үүсгэлт} -- тэмдэглэлийн агуулгад суурилсан зураг үүсгэх
  \item \textbf{Calendar view} -- стикерийг өдөртэй холбон харуулах
  \item \textbf{Push notification} -- тухайн өдрийн сануулга автоматаар илгээх
\end{itemize}

\vspace{0.3cm}
\noindent\textit{This study will be conducted at New Mongol College of Technology.}
